% Part: lambda-calculus
% Chapter: introduction
% Section: basic-pr-lambda

\documentclass[../../../include/open-logic-section]{subfiles}

\begin{document}

\olfileid[tr]{lam}{int}{bas}
\olsection{Temel İlkel Özyinelemeli Fonksiyonlar \usetoken{S}{lambda definable}}

\begin{lem}
$\Zero$, $\Succ$ ve $\Proj{n}{i}$ fonksiyonları !!{lambda definable}.
\end{lem}

\begin{proof}
$\Zero$, yalnızca $\lambd[x][\lambd[y][y]]$'dir.

Ardıl fonksiyonu $\Succ$, $\fn{Succ}(u)=\lambd[x][\lambd[y][x(uxy)]]$
ile tanımlanır. Bunun neden çalıştığını düşünmelisiniz: bir yineleyici olarak
ele alınan her $\num{n}$ sayısı ve her $f$ fonksiyonu için
$\fn{Succ}(\num{n},f)$, $y$ girdisinde $y$'den başlayarak $f$'yi $n$ kez,
ardından bir kez daha uygulayan bir fonksiyondur.

İzdüşümler hakkında söylenecek bir şey yoktur:
$\fn{Proj}^n_i(x_0,\dots,x_{n-1})=x_i$. Başka bir deyişle uzlaşımlarımıza
göre $\fn{Proj}^n_i$, $\lambd[x_0][\dots\lambd[x_{n-1}][x_i]]$ lambda
terimidir.
\end{proof}

\end{document}
